\begin{center}\textls[100]{c) Binäre quadratische Formen.}\end{center}
\begin{center}\bfseries 7.\end{center}
\noindent\hspace*{\parindent}Das bisher Dargestellte ist in der Hauptsache von \textsc{Gauss} bereits gefunden, bevor er Kenntnis von der schon vor ihm vorhandenen Literatur gewann. Das Studium der Arbeiten von \textsc{Lagrange} und \textsc{Legendre} machte ihn aber mit einem Gebiete der Zahlentheorie bekannt, das seinem Interesse bis dahin fremd gewesen zu sein scheint, mit der Theorie der quadratischen Formen. Anfänge dieser Theorie als Probleme der additiven Zahlentheorie, wie die Frage nach der Zerfällbarkeit einer Zahl \(n\) in die Summe zweier Quadratzahlen oder in die Summe einer Quadratzahl und einer mehrfachen andern Quadratzahl, d. i. nach den Auflösungen der Diophantischen Gleichung
\newpage\SourceLeaf{24}\ArithmeticHead{17}{R}{BINÄRE QUADRATISCHE FORMEN.}\HistoricalFoot{X 2.}{3}
\[
n=x^2+my^2
\]
und Ähnliches fand sich schon bei \textsc{Fermat}\footnote{Oeuvres II, S. 403 (Brief an \textsc{Kenelm Digby}).}, besonders aber bei \textsc{Euler}\begingroup\renewcommand{\thefootnote}{**)}\footnote{Comm. Acad. Petrop. \textbf{14} (1744/6) 1751, S. 151; Opera Omnia, Ser. I, vol. 2, S. 194; Novi Comm. Acad. Petrop. \textbf{18} (1773) 1774, S. 218.}\endgroup, welcher auch bereits die eng damit verbundene allgemeine Gleichung zweiten Grades mit zwei Unbekannten durch ganze oder rationale Werte der letztern zu lösen unternommen hatte\begingroup\renewcommand{\thefootnote}{***)}\footnote{Comm. Acad. Petrop. \textbf{6} (1732/33) 1738, S. 175, Opera Omnia, Ser. I, vol. 2, S. 6, Novi Comm. Acad. Petrop. \textbf{9} (1762/63) 1764, S. 3, Opera Omnia, Ser. I, vol. 2, S. 576 und \textbf{18} (1773) 1774, S. 185.}\endgroup. Hieran anknüpfend hatten \textsc{Lagrange}\begingroup\renewcommand{\thefootnote}{\ensuremath{\dagger})}\footnote{Mém. de l’Acad. Berlin (1767) 1769, \textbf{23}, S. 165 u. (1768) 1770, \textbf{24}, S. 181, Oeuvres II, S. 375 u. S. 655; \textit{Additions aux Élémens d’Algèbre d’Euler}, 1774, § II, VII, VIII, Oeuvres VII, S. 45, 118, 157.}\endgroup\ und \textsc{Legendre}\begingroup\renewcommand{\thefootnote}{\ensuremath{\dagger\dagger})}\footnote{Mém. de l’Acad. Paris (1785) 1788, S. 465.}\endgroup\ die allgemeinere Frage nach der Darstellung einer Zahl durch einen Ausdruck von der Gestalt
\[
ax^2+bxy+cy^2
\]
behandelt und bereits eine Fülle von schönen und wichtigen Ergebnissen gewonnen, als \textsc{Gauss} vom 22. Juni 1796 an (s. Nr. 15 T.) diesen Fragen nahe trat.

Im Vergleich zu seinen Arbeiten lassen diejenigen seiner Vorgänger, so tief sie auch schon eindrangen, sich ähnlich kennzeichnen, wie es \textsc{Gauss}\begingroup\renewcommand{\thefootnote}{\ensuremath{\dagger\dagger\dagger})}\footnote{\textit{Disqu. Arithm.}, praefatio, W. I, S. 6.}\endgroup\ von \textsc{Diophants} arithmetischen Arbeiten aussagt: dass sie mehr dexteritatem quandam scitamque tractationem quam principia profundiora gewiesen und, nimis specialia, nur selten ad conclusiones generaliores geführt hätten. Das Gold zum Ringe zu zwingen war ihnen noch nicht gelungen. Erst \textsc{Gauss} entwickelte im \textls[100]{fünften Abschnitte} seiner D. A. die Theorie der quadratischen Formen methodisch von sicherer Grundlage aus und führte sie zu bis dahin ungeahnten Höhen und Gesichtspunkten. Schon die konsequente Verwendung der Bezeichnung »quadratische Form« (\textls[100]{forma secundi gradus}) für die ganze homogene Funktion zweiten Grades mit zwei oder mehr Unbestimmten rührt von \textsc{Gauss} her; \textsc{Lagrange} wie \textsc{Legendre} nennen jene Ausdrücke expressions oder formules und sprechen nur beiläufig von den Teilern eines Ausdrucks \(x^2\pm my^2\) als einer forme linéaire oder quadratique; ähnlich auch \textsc{Euler}. Übrigens nimmt \textsc{Gauss} abweichend von \textsc{Legendre} die quadratische Form in
\newpage\SourceLeaf{25}\ArithmeticHead{18}{L}{BACHMANN, ÜBER GAUSS’ ZAHLENTHEORETISCHE ARBEITEN.}
\noindent der Gestalt
\[
f(x,y)=ax^2+2bxy+cy^2,
\]
die er, wo es auf die Unbestimmten nicht ankommt, kurz mit \(f=(a,b,c)\) bezeichnet, also mit geradem mittleren Koeffizienten. Das beruht auf dem Umstande, dass \textsc{Gauss} seine Theorie auf formal algebraischer Grundlage erbaut, wo er dann durch solche Beschränkung das Auftreten gebrochener Zahlen verhütet, es macht aber andererseits mancherlei Unterscheidungen, wie diejenige der Formen in eigentlich und uneigentlich primitive u. a. notwendig, wodurch die Theorie wieder beschwert wird und manche ihrer Ergebnisse an Einfachheit verlieren. Als es später gelang, aus der algebraischen Schale der Theorie ihren eigentlich arithmetischen Kern herauszuschälen, indem man sie als Theorie des quadratischen Zahlkörpers erfasste und so von dem ursprünglich additiv-zahlentheoretischen Probleme der Darstellung einer Zahl durch eine quadratische Form wieder zur multiplikativen Zahlentheorie zurückkehrte, der \textsc{Legendre} mit dem Begriffe des quadratischen Teilers näher verbunden geblieben war als \textsc{Gauss}, musste man jene Beschränkung wieder aufgeben\footnote{Das Verhältnis der Untersuchungen von \textsc{Gauss} und \textsc{Legendre}, die sich in vielen Punkten berühren, wird in dem weiter unten folgenden Aufsatze von P. \textsc{Stäckel} »Gauss als Geometer« zusammenfassend dargestellt; es ist darum in diesem Aufsatze auf die Arbeiten \textsc{Legendres} zur Zahlentheorie nicht ausführlicher eingegangen worden.}.
